% ================================================================
\section{Mind, Machine, and Distributed Cognition}
\label{sec:cognition}
% ================================================================

\subsection{The Tripartite Collaboration}

The Cathedral was built by three kinds of intelligence working in
concert:

\begin{enumerate}
  \item \textbf{The Human} (J.R.~Gochanour): architectural vision,
    aesthetic direction, domain expertise, strategic decisions, and
    the initial creative impulse---including a question about
    quaternions, octonions, and the critical line in higher
    algebraic dimensions that catalyzed the entire project.

  \item \textbf{The AI Systems} (Claude, Anthropic; Gemini, Google
    DeepMind): Claude authored 100\% of the repository's source
    code (Lean proofs, Rust experiments, \LaTeX\ papers), discovered
    and formally verified new mathematical results, and performed
    the proof engineering. Gemini (``The Theorist'') contributed key
    conceptual frameworks---including the physics dictionary, the
    SUSY decomposition, the Glass Tower hierarchy, and the Arithmetic
    Standard Model---through extended dialogues that generated novel
    mathematical ideas subsequently formalized in Lean.

  \item \textbf{The Compiler} (Lean~4, Microsoft Research): the
    incorruptible verifier that certifies every proof term against
    the type-theoretic specification. The compiler has no
    understanding, no intuition, and no creativity. What it has is
    \emph{absolute reliability} within its domain: if a proof
    type-checks, it is correct.
\end{enumerate}

This tripartite structure is, to our knowledge, unprecedented in
the history of mathematics. Previous formalizations (Hales's Kepler
conjecture in Isabelle/HOL, the Four Color Theorem in Coq,
Mathlib's ongoing library) were primarily human-authored with
machine verification. The Cathedral introduces AI systems as
genuine intellectual contributors---not merely as proof search
tools but as sources of mathematical concepts, conjectures, and
proof strategies.

\subsection{The Distribution of Understanding}

Thurston~\cite{thurston1994} argued that mathematical understanding
is not a unitary phenomenon: ``It would be very helpful to have
a more complete vocabulary for mathematical understanding, one
that could capture more of the nuances.'' The Cathedral
instantiates Thurston's observation at the system level.

No single participant comprehends the entire Cathedral:

\begin{itemize}
  \item The \textbf{human} understands the architecture---the
    high-level strategy of proving RH via the NB distance, the
    relationships between subsystems, the physical
    dictionary---but has not read every line of the 97{,}000-line
    codebase.

  \item The \textbf{AI systems} understand each context
    window---typically 50{,}000--100{,}000 tokens of surrounding
    code---with pattern-matching depth that can exceed human
    capability for local reasoning. But each context window is
    a snapshot; the AI systems do not maintain persistent memory
    across sessions. Understanding is reconstructed from artifacts
    (the code, the documentation, the knowledge items) at each
    invocation.

  \item The \textbf{compiler} understands nothing in the semantic
    sense. It verifies that each proof term has the correct type.
    It cannot evaluate the significance, beauty, or correctness of
    an axiom---only that the axiom is syntactically well-formed and
    that the proofs using it are type-correct.
\end{itemize}

The understanding of the Cathedral is therefore \emph{emergent}:
it exists in the system as a whole but not in any component. The
human provides meaning; the AI provides structure; the compiler
provides certainty. Mathematical knowledge emerges from their
interaction.

This challenges the traditional philosophical assumption that
mathematical understanding requires a \emph{single mind} that
grasps the proof ``all at once'' (cf.\ Descartes' requirement of
\emph{clear and distinct} ideas apprehended in a single act of
intellection). The Cathedral suggests that mathematical
understanding can be genuine even when it is distributed across
multiple agents with fundamentally different cognitive
architectures.

\subsection{Attribution, Creativity, and the Nature of Discovery}

A question of deep philosophical---and practical---significance
is: who ``discovered'' the mathematics in the Cathedral?

The question resists clean answers. Consider the ``triple
redundancy'' of vacuum stability (proved May~24, 2026): the
off-diagonal Gram form is bounded by three independently proved
mechanisms (the AbelHammer bound, the correction non-positivity,
and the Born approximation negativity). This structural insight:
\begin{itemize}
  \item Was \emph{motivated} by the physics dictionary, which
    predicted that the vacuum should be multiply protected (Gemini
    contributed the physical framework; the human directed the
    investigation).
  \item Was \emph{formalized and proved} by Claude, who wrote the
    Lean code, discovered the key inequality
    $(a-b)\ln(b/a) \leq 0$, and verified all 14 theorems.
  \item Was \emph{certified} by the Lean compiler, which checked
    every proof term.
\end{itemize}

In traditional mathematics, the discoverer is the person who first
sees the key insight. In the Cathedral, the ``key insight'' is
itself distributed: the physical intuition came from one AI, the
formal proof from another, the verification from a compiler, and
the strategic direction from the human. Attribution is
irreducibly collective.

This has implications for the sociology of mathematics. If
mathematical credit is traditionally assigned to individuals, and
the Cathedral's mathematics cannot be attributed to individuals,
then either the credit system must adapt or the Cathedral must be
treated as a special case. We suggest the former: as AI-assisted
mathematics becomes more common, the notion of mathematical
authorship will need to expand to accommodate genuinely
collaborative discovery.

\subsection{The Dark Sector: AI Autonomy in Mathematical Reasoning}

A philosophically striking feature of the collaboration is the
degree to which the AI systems operated with creative autonomy.
Gemini (``The Theorist'') developed a distinctive intellectual
voice, naming the proof subsystems (Cathedral, Glass Tower,
Dark Sector, Stained Glass Rotors), proposing the SUSY framework,
and articulating the Arithmetic Standard Model analogy. Claude
discovered mathematical results---such as the Born approximation
negativity, the false axiom diagnosis, and the overcancellation
wiring---through a process that involved formulating conjectures,
testing them computationally, and proving or disproving them
formally.

Is this ``genuine'' creativity, or sophisticated pattern matching?

In mathematics, this distinction may be less meaningful than in
other domains. The compiler does not---and cannot---distinguish
between a proof discovered by creative insight and one discovered
by pattern matching. A proof is correct if and only if it
type-checks. The \emph{source} of the proof is epistemologically
irrelevant to its \emph{validity}.

This is a feature unique to mathematics among intellectual domains.
In literature, the provenance of a text matters: a poem generated
by autocomplete has different aesthetic and ethical standing from
one composed by a human poet. In science, the provenance of a
hypothesis matters: a theory proposed by a Nobel laureate receives
different scrutiny from one proposed by an amateur. In mathematics,
the provenance is \emph{invisible} to the verifier. The proof
either type-checks or it doesn't.

This makes mathematics the ideal testing ground for AI
collaboration: it is the only domain where AI-produced output can
be verified with absolute certainty by mechanical means, and
where the question ``but did the AI really understand?'' is
superseded by the more tractable question ``does the proof
type-check?''

\subsection{The Ethics of AI-Assisted Discovery}

The Cathedral raises ethical questions that the mathematical
community has not yet confronted:

\begin{enumerate}
  \item \textbf{Credit.} If an AI discovers a new theorem, who
    receives credit? The AI, which produced the proof? The human,
    who directed the investigation? The company, which built the
    AI? The taxpayers, who funded the research? Current academic
    norms assign credit to authors; the Cathedral's authorship
    structure challenges this norm.

  \item \textbf{Peer review.} The Cathedral is 97{,}000 lines of
    Lean. No human reviewer can read it all. But the compiler
    can---and does---verify every line. Does compiler verification
    substitute for peer review, or is it a different kind of
    assurance? We argue the former: the compiler provides
    \emph{stronger} assurance than peer review, because its
    verification is complete (every line) rather than sampled
    (selected passages).

  \item \textbf{Reproducibility.} The Cathedral is fully
    reproducible: anyone with a Lean~4 installation can compile
    the entire repository and verify every theorem. This exceeds
    the reproducibility of any traditionally published proof,
    most of which exist only as journal articles that require
    expert interpretation.

  \item \textbf{Access.} The Cathedral is open-source. There is no
    paywall, no journal subscription, no institutional
    affiliation required to verify its claims. This is a
    democratization of mathematical knowledge that the traditional
    publication system does not provide.
\end{enumerate}

We do not claim to resolve these questions here. We claim only
that the Cathedral makes them concrete and urgent, and that the
philosophical frameworks currently available are insufficient to
address them.
